\section{Related Work}
\label{sec:related}

\subsection{EEG foundation models}

Large pretrained EEG encoders have proliferated in the past two years. LaBraM
(ICLR 2024 Spotlight)~\cite{labram2024} pretrains a $\sim$5--10\,M-parameter
transformer on 2{,}500 hours of EEG across 20 datasets with a vector-quantized
neural spectrum prediction objective, reporting state-of-the-art results on
TUH abnormal detection and several BCI tasks. EEGPT~\cite{eegpt2024} (NeurIPS
2024) reaches similar scale with a mask-based dual self-supervised objective
combining temporal and spatial alignment. NeuroLM~\cite{neurolm2025} (ICLR
2025) frames EEG as a foreign language with a text-aligned neural tokenizer
and language-model backbone, enabling multi-task learning. BIOT~\cite{biot2023}
generalizes the approach across biosignals (EEG, ECG, EMG) via power-spectral
tokenization. All of these are trained on large multi-GPU clusters with
training budgets measured in thousands of GPU-hours, and parameter counts
above 5\,M. Our work targets the substantially smaller scale where the entire
pretraining$+$evaluation pipeline runs on a single consumer GPU in under 15
minutes.

\subsection{LeJEPA and SIGReg}

LeJEPA~\cite{lejepa2025} introduced a theoretically-grounded
Joint-Embedding Predictive Architecture (JEPA) training objective that
combines next-embedding prediction with Sketched Isotropic Gaussian
Regularization (SIGReg). SIGReg uses the Cram\'er--Wold theorem
\cite{cramerwold1936}: enforcing one-dimensional Gaussianity on $M$ random
projections of an embedding distribution implies joint Gaussianity. The
resulting objective eliminates the stop-gradient, momentum encoder, and
teacher-student heuristics common to BYOL/DINO~\cite{byol2020,dino2021}, has
linear time and memory complexity, and achieves competitive performance on
ImageNet linear-probe with ViT-H/14 (79\,\%). LeWorldModel~\cite{lewm2026}
applied the LeJEPA recipe to a 15\,M-parameter world model trained on a
single GPU in hours, predicting frame embeddings 48$\times$ faster than
foundation-scale baselines while remaining competitive on 2D and 3D control
tasks. Our work brings the same compact, principled SIGReg-based recipe to
EEG, with a focus on small-corpus pretraining and edge-deployment.

\subsection{Concurrent work: Laya}

Independent and contemporaneous to ours, Panchavati et
al.~\cite{laya2026} (arXiv:2603.16281\,v2, March 2026) introduced Laya, an
application of LeJEPA to EEG foundation modelling at substantially larger
scale: 29{,}109 hours of EEG across 20{,}940 subjects and 17 channel
topologies drawn from TUH-EEG, HBN, NMT, MOABB, and EEGDash. Both reported
variants (Laya-S, Laya-B) use an effective batch size of 512 on two L40S
GPUs, learning rate $10^{-4}$ with cosine warm-up over 1{,}000 steps, weight
decay 0.05, and bf16 mixed precision. Laya-S trains for 10{,}000 steps on
10\,\% of pretraining data (3{,}000 hours); Laya-B trains for 20{,}000 steps
on the full corpus. Laya extends the LeJEPA objective with a Dynamic Channel
Mixer (queries-and-cross-attention with electrode-coordinate Fourier
features) plus a Query Specialization Loss. Architecturally, the resulting
recipe targets cross-dataset generalization at a substantially larger
compute budget than the regime we study (Laya-B uses
$\sim$$10\times 10^6$ effective sample-views, $\sim$16$\times$ ours).

Laya's downstream evaluation uses the multi-task EEG-Bench
framework~\cite{eegbench2025}, where each task is a multi-dataset benchmark
constructed by pooling several corpora. In particular, the
``LH vs RH MI'' task pools BCI Competition~IV-2a/2b, PhysioNet Motor Imagery
(EEGMMIDB), Weibo2014, Cho2017, Liu2022, Schirrmeister2017, Zhou2016, and
Kaya2018; the ``4-Class MI'' task pools BCI~IV-2a and Kaya2018. EEGMMIDB and
the BCI Competition IV datasets are explicitly excluded from Laya's
pretraining corpus to maintain the cross-dataset evaluation
property~\cite[Sec.~4]{laya2026}. Reported numbers
\cite[Tab.~1]{laya2026} are means across these multi-dataset pools
(Laya-S/Laya-B LH-vs-RH MI: 0.510 / 0.507 balanced accuracy; 4-Class MI:
0.273 / 0.288). On clinical tasks, Laya-S achieves the highest mean (0.603)
on the 10-task EEG-Bench clinical subset, with notable strength on TUH
abnormal detection (0.798), seizure detection (0.786), and mTBI
classification (0.699--0.815).

Our work is complementary in five concrete ways. First, where Laya pretrains
at foundation scale across many datasets and evaluates on cross-dataset
benchmarks, we pretrain on a single corpus (EEGMMIDB, $\sim$10\,h
$\sim$50 subjects) and evaluate in-distribution on the same corpus---the
operating regime realistic for individual research labs deploying compact
on-device EEG analytics. These are different research questions, both of
practical importance. Second, we explicitly characterize the scaling
behaviour of the small regime across six operating points (50--3{,}000
subjects equivalent, 200--30{,}000 training steps); Laya reports two
checkpoints and explicitly cites the absence of a scaling-law
characterization as a limitation~\cite[Sec.~5]{laya2026}. Third, we present
a SIGReg-weight ablation across $\lambda\in\{0, 0.1, 0.3, 1.0, 3.0, 10.0\}$
with full curves; Laya states the collapse threshold ($\lambda<0.02$)
without sweeping. Fourth, where Laya addresses channel heterogeneity via a
Dynamic Channel Mixer architecture, we ablate a simpler alternative---
controlled comparison between channel-padded cross-dataset transfer and
channel-matched pretraining---and reach a counterintuitive practitioner
conclusion. Fifth, we compare directly against supervised-from-scratch
training of the same architecture; Laya only benchmarks against other SSL
foundation models. Where both methods evaluate on EEGMMIDB-related motor
imagery, the comparison is necessarily indirect (Laya's number is a
multi-dataset pooled average, ours is single-corpus in-distribution); we
neither claim nor require Laya to under-perform our setup.

\subsection{Supervised baselines and motor-imagery BCI}

The classical BCI motor-imagery pipeline is filter-bank Common Spatial
Pattern features followed by an LDA or SVM classifier (FBCSP~\cite{fbcsp2008}),
which remains a strong supervised baseline at $\sim$65--75\,\% accuracy on
BCI Competition IV Dataset 2a in within-subject evaluation
\cite{eegnet2018,deepconv2017}. Deep-learning supervised baselines such as
EEGNet~\cite{eegnet2018} and ShallowConvNet~\cite{deepconv2017} reach
$\sim$70--75\,\% under the same protocol with $<$5{,}000 parameters. These
methods require per-task supervised training and do not amortize across
downstream tasks. Our SSL$+$linear-probe pipeline complements rather than
replaces them: a single pretrained encoder serves multiple downstream tasks
at sub-second probe cost, with absolute accuracy slightly below FBCSP/EEGNet
but with reusability and edge-deployability advantages.

\subsection{Anti-collapse SSL beyond SIGReg}

Alternative anti-collapse mechanisms include the explicit
variance/invariance/covariance term of VICReg~\cite{vicreg2022}, the
redundancy reduction of Barlow Twins~\cite{barlowtwins2021}, the
momentum-encoder of BYOL~\cite{byol2020} and DINO~\cite{dino2021}, the
asymmetric stop-gradient of SimSiam~\cite{simsiam2021}, and the contrastive
formulation of SimCLR/MoCo~\cite{simclr2020,moco2020}. SIGReg replaces these
mechanisms with a single principled distribution-matching loss, eliminating
the multiple hyperparameter schedules that plague competitive methods. Our
$\lambda$ ablation (\cref{sec:results}) provides further empirical
support for SIGReg's stability claim: a 100$\times$ change in $\lambda$
changes downstream accuracy by 10\,pp but never destabilizes training.
